% SOURCE_AUTHORITY: sga5_fr_workpass.tex, lines 15068--15074 (LF-normalized unit).
% SOURCE_SHA256: 791F4EFFC5E02832D5D77ED03518C8156D6F07E4C8238B03545DB93D883FBB28
Utilizemos agora os lemas anteriores para nos reduzirmos ao caso em que $X$ seja a reta projetiva.

1\textsuperscript{a} redução: por indução noetheriana pode-se construir uma partição $(U_j)_{1\le j\le r}$ de $X$ por esquemas afins tal que $U_{j+1}$ seja aberto em $X-\bigcup_{i=1}^jU_i$ para $j=1,\ldots,r-1$. Suponhamos estabelecido o teorema para todo par $(U,G)$ formado por um esquema afim $U$ de tipo finito sobre $\mathbb F_p$ e um $\Omega$-feixe construtível $G$ sobre $U$; é, portanto, verdadeiro para $(U_j,F|U_j)$, $j=1,\ldots,r$; graças ao lema 3a), deduz-se por indução descendente sobre $j$ que é verdadeiro para $(\bigcup_{i\ge j}U_i,F|\bigcup_{i\ge j}U_i)$, e, portanto, para $(X,F)$, já que $X=\bigcup_{i\ge1}U_i$. Basta, então, provar o teorema quando $X$ é afim.

2\textsuperscript{a} redução: suponhamos $X$ afim, de dimensão $n$; existe um morfismo finito $u:X\to\mathbb E^n$ de $X$ ao espaço afim de dimensão $n$. Raciocinemos por indução sobre $n$. Para $n=0$, $X$ é finito sobre $e$ e o teorema é verdadeiro pelo corolário do lema 1. Se $n\ge1$, suponhamos estabelecido o teorema para os esquemas afins de dimensão $n-1$; sejam $u:X\to\mathbb E^n$ um morfismo finito e $v:\mathbb E^n\to\mathbb E^1$ um morfismo de projeção, definido a partir de um isomorfismo $\mathbb E^n\simeq\mathbb E^1\times\mathbb E^{n-1}$; ao compô-los obtemos um morfismo $w=v\circ u:X\to\mathbb E^1$ cujas fibras são esquemas afins de dimensão $n-1$. O lema 1 mostra então que a propriedade $(P_{F,w})$ é verdadeira pela hipótese de indução. Graças ao lema 2, vê-se que basta estabelecer o teorema para os pares $(\mathbb E^1,\R^i_!w(F))$. Assim, reduz-se ao caso em que $X=\mathbb E^1$ seja a reta afim.

3\textsuperscript{a} redução: suponhamos $X=\mathbb E^1$; seja $i:\mathbb E^1\to\mathbb P^1$ a imersão da reta afim na reta projetiva $\mathbb P^1$, e seja $Z=\mathbb P^1-i(\mathbb E^1)$ o subesquema fechado reduzido ao ponto no infinito de $\mathbb P^1$. Como o teorema é verdadeiro para $(Z,i_!(F)|Z)$, pelo lema 3a) equivale estabelecê-lo para $(\mathbb E^1,F)$ ou para $(\mathbb P^1,i_!(F))$ (como $i$ é uma imersão aberta, tem-se $F=i_!(F)|\mathbb E^1$).
